 % Første linje er dokumentklassen. Brug *altid* memoir!
 \documentclass[a4paper,oneside,article]{memoir}

 % Herunder indlæses forskellige pakker
 \usepackage[utf8]{inputenc}  % Korrekt håndtering af æ, ø og å
 \usepackage[T1]{fontenc}  % Korrekt håndtering af æ, ø og å
 \usepackage{microtype}  % Typografisk magi! Giver bl.a. pænere orddeling
 \usepackage{graphicx}  % Gør det muligt at indsætte billeder
 \usepackage{amsmath}  % Giver adgang til uundværlige matematikting
 \usepackage{mathtools}  % Retter ting i ams og giver ekstra muligheder
 \usepackage[danish]{babel}  % Danske betegnelser og orddeling
 \usepackage{lmodern}
 \usepackage{alphabeta}
 \usepackage{float} %Mere kontrol over billede-placering
 \renewcommand{\danishhyphenmins}{22}  % Bedre dansk orddeling
 \usepackage[exponent-product=\ensuremath{\cdot}]{siunitx} % Sientific notation for forskellige ting, med en regel om at den skal bruge cdot istedet for times. 
 \usepackage{derivative} % Til at skrive differentialligninger 
 \usepackage{xcolor} % For at få flere tekst farver 
 \usepackage[colorlinks=true]{hyperref} % For at få refrences man kan trykke på, slet [] hvis man vil have kasser i stedet for farvet tekst.
 \usepackage{pdfpages} % Til at indsætte andre pdf'er
 \usepackage{subfig} % For at have to billeder ved siden af hinanden
 \usepackage{unicode-math}
 \usepackage{bm}
 
 % Pænere toc indents 
 \usepackage{tocloft}
 \cftsetindents{section}{1em}{2em}
 \cftsetindents{subsection}{4em}{0em}
 
 
 % Indholdet i dit dokument starter her!
 \begin{document}
 \author{Mikkel Moth Billing \\ 202307989}
 \title{Kvantemekanik Formelsamling}
 \date{\today}
 \maketitle   % Indsætter overskriften
 % Table of Contents indstillinger
 \setcounter{tocdepth}{2}
 \setcounter{secnumdepth}{1}
 \hypersetup{linkcolor=black} % Laver tableofcontents sort.  
 \tableofcontents %Indsætter ToC
 \hypersetup{linkcolor=red} % og så alle andre hyperrefs røde. 
 
 %Custom Commands
 %Til opgavebeskrivelser
 \newcommand{\opg}[1]{\textcolor{gray}{\emph{#1}}\plainbreak{1}} 
 
 \newcommand{\ti}[1]{\cdot 10^{#1}} % Let: gange 10^x
 \newcommand{\e}[1]{\emph{e}^{#1}} % Pænere e^x
 \newcommand{\vm}[2]{\begin{pmatrix} #1 \\ #2 \end{pmatrix}}
 \newcommand{\mat}[1]{\begin{bmatrix} #1 \end{bmatrix}}
 
 % Lette paranteser 
 \newcommand{\br}[2][1]{%
     \ifcase#1
     \or \left( #2 \right) %1 eller ingenting
     \or \left[ #2 \right] %2
     \or \left\langle #2 \right\rangle %3
     \or \left\{ #2 \right\} %4
     \or \left| #2 \right| %5
     \fi
 }
 %Til at sætter bileder ind i store dokumenter
 \newcommand{\bil}[1]{
 \begin{figure}[H]
     \centering
     \includegraphics[width=0.75\linewidth]{Billeder/#1.png}
     \label{fig:#1}
 \end{figure}
 }
 
 %Til at sætte pdf'er ind i store dokumenter (kun pdf'er på mere end 1 sider)
 \newcommand{\pdf}[4]{
 \includepdf[pages=#3,scale=.85,pagecommand={\subsection{#2 - \ref{ch:opg}}  \label{asc:#2}}]{pdf/#1.pdf}
 \includepdf[pages=\the\numexpr 1 + #3\relax-#4, scale=.85]{pdf/#1.pdf}
 }
 
 \newcommand{\h}{\hbar}
 \newcommand{\dd}{\partial}
 \newcommand{\eq}[1]{\label{eq#1}\tag{#1}}
 \renewcommand{\a}{\hat{a}}
 \newcommand{\n}{\plainbreak{1}}
 \newcommand{\ep}{\epsilon}
 
 \newcommand{\hv}[2][]{%
     \if\relax\detokenize{#1}\relax
         |#2\rangle%
     \else
         \if\relax\detokenize{#2}\relax
             \langle #1|%
         \else
             \langle #1|#2\rangle%
         \fi
     \fi
 }
 
 
 
 % Lav ny pagestyle (kun odd fordi enkeltsidet)
 \makepagestyle{Title}
 \pagestyle{Title}
 %\makeoddhead{Title}{\footnotesize\color{gray}{\theauthor}}{}{\textcolor{gray}{\thedate}}
 \makeoddfoot{Title}{}{\footnotesize\color{gray}{\thepage{} / \thelastpage}}{}
 % Også fod på siden med \maketitle
 \makeoddfoot{plain}{}{\footnotesize\color{gray}{\thepage{} / \thelastpage}}{}
 \newpage
 
 %Skriv
 \chapter{Math}
 \section{Vectors}
\begin{align*}
    &a(\bar{A}+ \bar{B}) = a\bar{A} + a\bar{B} \\
    &\bar{A} \cdot \bar{B} = |\bar{A}||\bar{B}|\cos(\theta) \eq{1.1} \\    
    &\bar{A} \times \bar{B} = |\bar{A}||\bar{B}|\sin(\theta)\hat{n} \eq{1.4} \\
    &(\bar{A}\times \bar{B}) = -(\bar{B}\times \bar{A}) \eq{1.6} \\
    &\bar{A} \cdot (\bar{B} \times \bar{C}) = \bar{B} \cdot (\bar{C} \times \bar{A}) = \bar{C} \cdot (\bar{A} \times \bar{B}) \eq{1.15} \\
    &\bar{A} \times (\bar{B} \times \bar{C}) = \bar{B}(\bar{A} \cdot \bar{C}) - \bar{C}(\bar{A} \cdot \bar{B}) \eq{1.17} \\
    &(\bar{A} \times \bar{B}) \cdot (\bar{C}\times \bar{D}) = (\bar{A} \cdot \bar{C})(\bar{B} \cdot \bar{D}) - (\bar{A} \cdot \bar{D})(\bar{B} \cdot \bar{C})  \\
    &\bar{A} \times [\bar{B} \times (\bar{C} \times \bar{D})] = \bar{B}[\bar{A}\cdot(\bar{C}\times \bar{D})] - (\bar{A}\cdot\bar{B})(\bar{C}\cdot\bar{D})\eq{1.18}
\end{align*} 

\section{Calculus}
\begin{align*}
    dT &= \br{\frac{\dd T}{\dd x}+\hat{x}+\frac{\dd T}{\dd y}\hat{y}+\frac{\dd T}{\dd z}\hat{z}} \cdot (dx \hat{x} + dy \hat{y} +dz \hat{z}) \\
    &= (\nabla T) \cdot d\bar{l}, \eq{1.35} \\
\end{align*}
where the gradent ($\nabla$) is an operator: 
\begin{align*}
    \nabla = \hat{x}\frac{\dd}{\dd x} + \hat{y}\frac{\dd}{\dd y} + \hat{z}\frac{\dd}{\dd z}. \eq{1.39}
\end{align*}

The gradient is used to calculate the divergence of a vector with: 
\begin{align*}
    \nabla \cdot \bar{v} = \frac{\dd v_x}{\dd x} + \frac{\dd v_y}{\dd y} + \frac{\dd v_z}{\dd z}, \eq{1.40}
\end{align*}
and the curl with: 
\begin{align*}
    &\nabla \times \bar{v} = \begin{vmatrix} \hat{x} & \hat{y} & \hat{z} \\ \frac{\dd}{\dd x} & \frac{\dd}{\dd y} & \frac{\dd}{\dd z} \\ v_x & v_y & v_z \end{vmatrix} \\ 
    = \left(\frac{\dd v_z}{\dd y} - \frac{\dd v_y}{\dd z}\right)\hat{x} &+ \left(\frac{\dd v_x}{\dd z} - \frac{\dd v_z}{\dd x}\right)\hat{y} + \left(\frac{\dd v_y}{\dd x} - \frac{\dd v_x}{\dd y}\right)\hat{z}. \eq{1.41}
\end{align*}

The product rules for the gradient are: 
\begin{align*}
    &\nabla (fg) = f\nabla g + g\nabla f \\
    &\nabla (\bar{A} \cdot \bar{B}) = \bar{A} \times (\nabla \times \bar{B}) + \bar{B} \times (\nabla \times \bar{A}) + (\bar{A} \cdot \nabla)\bar{B} + (\bar{B} \cdot \nabla)\bar{A} \\
    &\nabla \cdot (f\bar{A}) = f(\nabla \cdot \bar{A}) + \bar{A} \cdot (\nabla f) \\
    &\nabla\cdot(\bar{A}\times \bar{B}) = \bar{B}\cdot(\nabla\times \bar{A}) - \bar{A}\cdot(\nabla\times \bar{B}) \\
    &\nabla\times(f\bar{A}) = f(\nabla\times \bar{A}) + \nabla f \times \bar{A} \\
    &\nabla\times(\bar{A}\times \bar{B}) = \bar{A}(\nabla\cdot \bar{B}) - \bar{B}(\nabla\cdot \bar{A}) + (\bar{B}\cdot \nabla)\bar{A} - (\bar{A}\cdot \nabla)\bar{B}
\end{align*}

When we use the gradient to make a second derivitive we get the Laplacian: 
\begin{align*}
    \nabla (\nabla T) = \frac{\dd^2 T}{\dd x^2} + \frac{\dd^2 T}{\dd y^2} + \frac{\dd^2 T}{\dd z^2}, \eq{1.42}
\end{align*}
and the laplacian of a vector: 
\begin{align*}
    \nabla^2\bar{v} = (\nabla^2 v_x)\hat{x} + (\nabla^2 v_y)\hat{y} + (\nabla^2 v_z)\hat{z}. \eq{1.43}
\end{align*}

We also have the following identities:
\begin{align*}
    &\nabla \times (\nabla T) = 0 \eq{1.45}\\
    &\nabla \cdot (\nabla \times \bar{v}) = 0 \eq{1.46}\\
    &\nabla \times (\nabla \times \bar{v}) = \nabla(\nabla \cdot \bar{A}) - \nabla^2 \bar{A}. \eq{1.47}
\end{align*}

\section{Fundamental theorem of calculus}
Gradient: 
\begin{align*}
    \int_{\bar{a}}^{\bar{b}} (\nabla T) \cdot d\bar{l} = T(\bar{b}) - T(\bar{a}) \eq{1.55} 
\end{align*}
Divergence:
\begin{align*}
    \int_{\mathcal{V}} (\nabla \cdot \bar{v})dV = \oint_{\mathcal{S}} \bar{v} \cdot d\bar{a} \eq{1.56}
\end{align*}
Curl: 
\begin{align*}
    \int_{\mathcal{S}} (\nabla \times \bar{v}) \cdot d\bar{a} = \oint_{\mathcal{P}} \bar{v} \cdot d\bar{l} \eq{1.57}
\end{align*}

\chapter{Electrostatics}
\section{Electric field}
The force on a test charge $(Q)$ from a single point charge $(q)$, a distance $\bar{r}$ away, is calculated by Coulomb's law: 
\begin{align*}
    \bar{F} = \frac{1}{4 \pi \ep_0}\frac{q Q}{r^2}\hat{r}, \eq{2.1}
\end{align*}
where  $\ep_0 = 8.85\ti{-12}\frac{\text{C}^2}{\text{N}\cdot \text{m}^2}$ is the permittivity of free space. 


\end{document}